\documentclass[11pt,a4paper]{article}
\usepackage[margin=2.5cm]{geometry}
\usepackage[T1]{fontenc}
\usepackage[utf8]{inputenc}
% amssymb supplies \gtrsim (used when describing cluster charge thresholds).
% Without it pdflatex aborts on an undefined control sequence -- which went
% unnoticed because this guide had never actually been compiled.
\usepackage{amsmath}
\usepackage{amssymb}
% The Live tab's workflow buttons are labelled (1)..(5) in the GUI, and this
% guide names them the same way, with the same circled digits. pdflatex's utf8
% layer has no glyph for U+2460..U+2464, so map them onto pifont's circled
% digits rather than rewriting every mention as plain text.
\usepackage{pifont}
\DeclareUnicodeCharacter{2460}{\ding{172}}
\DeclareUnicodeCharacter{2461}{\ding{173}}
\DeclareUnicodeCharacter{2462}{\ding{174}}
\DeclareUnicodeCharacter{2463}{\ding{175}}
\DeclareUnicodeCharacter{2464}{\ding{176}}
\usepackage{booktabs}
\usepackage{longtable}
\usepackage{graphicx}
\usepackage{xcolor}
\usepackage{hyperref}
\hypersetup{colorlinks=true, linkcolor=blue!60!black, urlcolor=blue!60!black}
\newcommand{\param}[1]{\texttt{\textbf{#1}}}
\newcommand{\gui}[1]{\textsf{#1}}

\title{qBOUNCE CMOS Neutron Detector\\\large Pipeline \& Live Acquisition --- User Guide}
\author{}
\date{\today}

\begin{document}
\maketitle
\tableofcontents

\section{What this software does}
A CMOS camera (IDS, $5472\times3672$ px, Mono8) watches a boron-10 converter
layer. A neutron capture, $n + {}^{10}\mathrm{B} \rightarrow \alpha +
{}^{7}\mathrm{Li}$, emits the two daughters \emph{back-to-back}: only one of
them reaches the sensor and leaves a bright, compact track ($\geq$ 5\,px).
The software:
\begin{enumerate}
  \item learns the per-pixel noise of the sensor from dark frames
        (\emph{background model}: mean, $\sigma$, dead-pixel map);
  \item detects clusters of pixels above $\mathrm{mean} + n_\sigma\,\sigma$;
  \item classifies each cluster with a machine-learning model
        (\texttt{lithium} = heavy track, \texttt{gamma}, \texttt{artifact});
  \item counts neutrons \emph{live}: $N_{\mathrm{neutrons}} =
        N_{\alpha} + N_{\mathrm{Li}}$ (one daughter per capture);
  \item estimates the deposited energy of each track
        (\param{total\_charge} $= \sum_{\mathrm{pixels}}(\mathrm{ADU} -
        \mathrm{background})$) and builds the $\alpha$/$^{7}$Li energy
        spectrum;
  \item monitors camera health (hot/dead pixels, radiation damage).
\end{enumerate}

\section{Installation on the offline Windows~10 machine}
\textbf{Step A — on any machine WITH internet} (e.g.\ the development Mac):
\begin{verbatim}
python3 PrepareSetupLab.py
\end{verbatim}
This builds \texttt{ExportLab/} — the complete offline kit: sources, GUI,
docs, the trained classifier, every Python wheel for Windows (including
\texttt{cmake} and \texttt{ninja} as pip wheels), the prebuilt OpenCV C++
package, and a generated \texttt{cpp\_requirements.txt} manifest (parsed from
\texttt{CMakeLists.txt}'s \texttt{find\_package} calls --- the closest thing
to a ``pipreqs for C++'').

\textbf{Step B — on the offline lab machine}: copy \texttt{ExportLab} to
e.g.\ \texttt{C:\textbackslash qBounce}. \texttt{WindowsLauncher.py} works
fully offline (venv + wheels, installs cmake/ninja into the venv with no
admin rights, extracts the bundled OpenCV and passes
\texttt{-DOpenCV\_DIR} to CMake, compiles the detector) and has three
deliberately separate usages, so that the slow ``provision the machine''
work is not repeated every single day:
\begin{itemize}
  \item \texttt{python WindowsLauncher.py --check} --- verifies the setup
        \textbf{only} ($\sim$10\,s): no build, no GUI, nothing changed on
        disk. Run it once on arrival, or whenever something looks wrong.
  \item \texttt{python WindowsLauncher.py} --- the \textbf{full run}: venv +
        packages + OpenCV + CMake build, then launches the GUI.
  \item \texttt{python WindowsLauncher.py --no-launch} --- everything above
        \emph{except} starting the GUI at the end (used by
        \texttt{LabValidation.py}, below, to provision unattended).
\end{itemize}
Run the full \texttt{WindowsLauncher.py} \textbf{once}. After that
provisioning has succeeded, the \textbf{daily} start no longer goes through
the launcher at all:
\begin{itemize}
  \item \gui{CheckSetup.bat} --- double-click to re-run \texttt{--check}
        (verification only, safe to run any time something looks wrong).
  \item \gui{RunLab.bat} --- double-click for the \textbf{daily} start: runs
        \texttt{Code\textbackslash venv\textbackslash Scripts\textbackslash
        python.exe PipelineController.py} directly from the
        already-provisioned venv, and refuses with a clear message if that
        venv does not exist yet (i.e.\ if \texttt{WindowsLauncher.py} was
        never run to completion).
\end{itemize}

\textbf{Validating the kit}: \texttt{python LabValidation.py} exercises the
whole chain unattended --- runs \texttt{WindowsLauncher.py --check}, then a
full \texttt{--no-launch} provisioning, proves the resulting
\texttt{detector(.exe)} actually \textbf{starts} (i.e.\ that the OpenCV DLLs
resolve), and finally replays the exact commands the GUI's \gui{Build}
button would issue. Every subprocess has a timeout, so a run can never hang
overnight waiting for input. Walk away and read
\texttt{lab\_validation\_report.txt}, written next to the script, for a
PASS/FAIL line per check; exit code 0 means every check the machine could
run, passed.

Two things the kit cannot bundle: the \textbf{C++ compiler} and the
\textbf{IDS peak SDK}.
\begin{itemize}
  \item \textbf{Compiler}: must be \textbf{MSVC} (Visual Studio Build
        Tools --- create an offline layout with
        \texttt{vs\_buildtools.exe --layout} on an online PC if missing; the
        launcher detects and explains this). \textbf{MinGW is not a
        substitute}: the bundled OpenCV binaries are a \texttt{vc16} build,
        and MinGW cannot link against it. A machine with MinGW on
        \texttt{PATH} would otherwise look like it ``has a compiler'' and
        only fail later, at the link step, which is why the launcher checks
        specifically for MSVC rather than accepting any \texttt{g++}.
  \item \textbf{IDS peak SDK}: still install it from its own installer for
        real-camera capture (drivers + GenTL producer). Its Python
        \emph{bindings}, however, are \textbf{not included in the standard
        installer} (verified on IDS peak 26.06, on a real machine) --- the
        kit no longer expects to find them under \texttt{Program Files}.
        Instead \texttt{PrepareSetupLab.py} downloads \texttt{ids\_peak} and
        \texttt{ids\_peak\_ipl} straight from PyPI into
        \texttt{Code\textbackslash wheels}, so the offline machine installs
        them like any other Python dependency. Only the Python layer travels
        with the kit --- the native SDK (drivers/GenTL) must still be
        installed on the lab machine itself.
\end{itemize}

\section{Standard workflow (Pipeline tab)}
Run the steps top to bottom. Each step's output feeds the next.
\begin{enumerate}
  \item \gui{Clean} --- deletes previous outputs (background model,
        detections CSVs, result plots). \emph{Keeps} the hand-labelled
        annotation sample and the trained ML model.
  \item \gui{Build} --- compiles the C++ detector (\texttt{cmake} + \texttt{make}).
  \item \gui{Background Model} --- runs the detector over the \emph{dark}
        folder. Produces \texttt{background\_model.yml} (binary format:
        per-pixel mean, variance accumulator, sample count, dead-pixel mask).
        \textbf{Redo this whenever exposure, gain, temperature or the camera
        changes} --- the model is specific to the acquisition conditions.
  \item \gui{Cluster Detection} --- runs the detector over the \emph{signal}
        folder (resuming from the background model) $\rightarrow$
        \texttt{detections.csv}: one row per cluster with 20 features
        (position, size, charge, shape, SNR\ldots).
  \item \gui{Labeling} --- opens the annotation GUI on a sample of clusters.
        You assign the true class of each cluster (hotkeys shown on the
        buttons; \texttt{[U]}/\texttt{Ctrl-Z} undoes a mistake). Because
        $\alpha$ and $^{7}$Li tracks cannot be told apart visually, both are
        labelled \texttt{lithium} (heavy track).
  \item \gui{Classify / Train} --- trains a Random-Forest on your labelled
        sample (quality = cross-validated F1, shown in the GUI) and exports
        it as \texttt{model\_bundle.joblib} + \texttt{model\_bundle.onnx}.
        \gui{Predict} then applies the model to the \emph{whole}
        \texttt{detections.csv} $\rightarrow$
        \texttt{detections\_classified.csv} + summary plots. This is the
        offline neutron count of the run.
  \item \gui{Backtest} --- re-runs the model on a \emph{golden dataset} (a
        CSV with a human \texttt{true\_label} column, chosen in
        \gui{Parameters} $\rightarrow$ \gui{Backtesting}) and reports
        precision/recall per class. Use an independently labelled sample,
        not the training sample, for an honest number.
  \item \gui{Live} tab --- real-time acquisition and neutron counting
        (Section~\ref{sec:live}).
\end{enumerate}

\section{Parameters reference (Parameters tab)}

\subsection{Acquisition presets (IDS camera)}
\begin{longtable}{p{4.2cm}p{10.5cm}}
\toprule
\textbf{Parameter} & \textbf{Meaning and impact} \\
\midrule
\param{Acquisition mode} &
\textbf{Energy resolution} (default): reduced exposure/gain so heavy tracks
stay below the 255 saturation level $\rightarrow$ the deposited-energy
spectrum is usable. \textbf{High sensitivity}: full exposure/gain (historical
behaviour); use when faint events (gammas) matter more than energy
resolution. \\
\param{Exposure us (energy)} &
Exposure time of the energy preset. Tune until \texttt{study\_energy}
reports \emph{Saturated} $\approx 0\%$. Shorter exposure = darker tracks,
fewer per frame; note that counting efficiency is the \emph{duty cycle}
$= \mathrm{exposure} \times \mathrm{fps}$ --- if you shorten the exposure,
raise the fps accordingly or you will miss neutrons arriving between
frames. \\
\param{Gain dB (energy)} & Analog gain of the energy preset. Start at 0. \\
\param{Target saturated px/frame} &
The only knob of \gui{Auto-tune exposure}: acceptable mean number of
saturated pixels ($\geq 254$) per frame. It is an \emph{absolute pixel
count}, not a percentage: with $\sim$40 tracks $\times$ $\sim$10 px each, a
target of 50 still allows $\sim$10\% of track area to clip; use 10--20 for
clean spectroscopy. \\
\param{Auto-tune exposure (beam ON)} &
Halves the exposure from 500\,ms until the target is met, then fills the
energy-mode exposure field. Run it with the beam ON --- it tunes on real
track brightness. If it hits the 100\,\textmu s floor, lower the gain and
re-run. \\
\bottomrule
\end{longtable}

\subsection{Detector parameters (C++)}
\begin{longtable}{p{4.2cm}p{10.5cm}}
\toprule
\textbf{Parameter} & \textbf{Meaning and impact} \\
\midrule
\param{n\_sigma} (default 5) &
Detection threshold: a pixel is an event if
$\mathrm{ADU} > \mathrm{mean} + n_\sigma \sigma$. Lower = more sensitive
and more noise clusters; higher = misses faint tracks.
\texttt{src/study\_nsigma.py} plots cluster statistics vs.\ $n_\sigma$ to
choose it objectively. \\
\param{warmup\_frames} &
Number of dark frames used to build the dead-pixel map (auto = the whole
dark folder). \\
\param{Min cluster size (px)} (default 4) &
Clusters smaller than this are \emph{dropped at the source} (no CSV row, no
classification). Heavy tracks are $\geq 5$ px; gammas/betas are 1--3 px.
Set to 0 to record everything again (e.g.\ to label gammas). Dropping junk
reduced $\sim$11\,000 clusters/frame to $\sim$60 on real data with no loss
of heavy tracks. \\
\param{Min cluster charge (ADU)} (default 300) &
Same idea on deposited charge. Heavy tracks carry $\gtrsim 700$ ADU;
noise/artifacts $\lesssim 70$. \\
\param{Cluster gap tolerance (px)} (default 0) &
Morphological closing radius that bridges small gaps in split tracks
(useful for muons). 0 = off. \\
\param{Cluster grow radius (px)} (default 1) &
After clustering, adds neighbouring pixels above \param{Grow min SNR} to the
footprint: recovers the sub-threshold halo charge around a track (better
energy proxy). \textbf{Changing it changes the charge scale} --- redo the
energy calibration afterwards. \\
\param{Grow min SNR} (default 1.0) & Minimum per-pixel SNR for grown pixels. \\
\bottomrule
\end{longtable}

Two protections are automatic (no knob): pixels saturated/black during most
of the warmup are masked, and any pixel that is an \emph{event for 5
consecutive frames} is declared hot and masked permanently (real particles
never hit the same pixel every frame) --- this catches radiation damage that
appears \emph{during} a run.

\subsection{Matched operating point (important)}
The background model, the cluster filters and the trained classifier form a
\textbf{matched set}, tied to one gain. If you lower the gain to de-saturate
tracks for the energy spectrum, three things move at once: the noise $\sigma$
(so \textbf{rebuild the background model} at the new gain), the absolute ADU
scale (so use \param{Min peak SNR}, which is gain-invariant, not
\param{Min cluster charge}), and the classifier's input distribution --- a
model trained on 100\%-saturated tracks has never seen a de-saturated one and
may mislabel it. Three safeguards make this visible: (i) the model records
the saturation it was trained at, and the live view pops a one-time
\emph{operating-point mismatch} warning if the live saturation diverges by
$>$25 points; (ii) \gui{Check separability} (Results tab) tests whether signal
and noise can still be told apart in gain-robust features and shows a red
dialog if they have MERGED --- at which point the neutron count is
meaningless; (iii) the pre-filters are a throughput guardrail only (the
classifier does the physics), set loose enough not to touch real tracks.
The clean way to run de-saturated: pick a gain (\gui{Auto-tune gain}), rebuild
the background model, take a labelling sample \emph{at that gain}, retrain,
and keep the gain fixed --- so the model always matches the saturation it was
trained on.

\subsection{Live parameters (Live tab)}\label{sec:live}
The \gui{Live Source} box only shows the fields relevant to the selected
source: \gui{Simulated folder}/\gui{Simulated FPS} appear for the replay
source and disappear for \gui{IDS camera} (whose settings live in
\gui{Parameters}). Everything below the source box --- \param{Duration},
batching, the model bundle, \gui{Start}/\gui{Stop} --- applies identically
to \emph{either} source: there is one acquisition run, one Start control
(the workflow's ②), regardless of where frames come from.

\emph{Command-line note}: the GUI always passes the source explicitly, but
if \texttt{LiveDetectionEngine.py} is ever invoked directly (e.g.\ for
debugging), \param{--camera \{simulated,ids\}} is \textbf{mandatory} --- it
has no default. This is deliberate: a stray invocation can never silently
fall back to simulated frames when a real measurement was intended.
\begin{longtable}{p{4.2cm}p{10.5cm}}
\toprule
\textbf{Parameter} & \textbf{Meaning and impact} \\
\midrule
\param{Duration (s)} & Total length of the live session (both sources); the
engine auto-stops after this many seconds. \\
\param{Batch seconds} / \param{Batch max frames} &
A \emph{batch} is a display/classification grouping of frames (whichever
limit is reached first closes the batch). It does not affect the physics ---
frames are streamed to the detector one by one regardless. With
\param{Batch max frames} = 1 you get one telemetry line per frame. \\
\param{Rolling window (s)} &
Time window of the rolling neutron rate. Poisson statistics: the relative
error shown ($\pm$\%) is $1/\sqrt{N}$ with $N$ = neutrons in the window;
$N \geq 100$ for $\sim$10\% precision. \\
\param{Model bundle} &
The trained classifier; the live engine actually loads the
\texttt{.onnx} + \texttt{.meta.json} exported next to it. \\
\param{Save during live} &
What is written to disk while measuring, each on its own writer thread so
detection is never throttled. \textbf{raw frames}: bit-identical
\texttt{.bmp} archive of every camera frame (IDS-Cockpit replacement,
$\sim$20\,MB/frame). \textbf{data\_analysed}: compact products only
($\sim$20\,KB/frame masked images + noise model + arrival map). Default:
both. \\
\param{Raw recording folder} & Destination of the raw \texttt{.bmp} archive. \\
\bottomrule
\end{longtable}

\paragraph{Storage budget (4\,TB server, 20.1\,MB/frame, 2\,fps, 24\,h/day).}
\begin{center}
\begin{tabular}{lrr}
\toprule
mode & disk per day & 4\,TB lasts \\
\midrule
raw frames            & $\sim$3.5\,TB & $\sim$1.2 days \\
data\_analysed only   & $\sim$4\,GB   & $\sim$1000 days \\
\bottomrule
\end{tabular}
\end{center}
\emph{Caveat}: data\_analysed keeps only the pixels of clusters that passed
the quality filters — it is an analysis product, not a raw archive. A
lower detection threshold cannot be re-applied to it afterwards; keep raw
frames (or lower the filters) whenever that flexibility matters.

\paragraph{Temporal coverage.} Next to the fps choice, a coverage label
shows $\mathrm{exposure} \times \mathrm{fps}$ — the fraction of wall-clock
time the sensor integrates. Neutrons arriving outside it are never counted;
divide any measured rate by the coverage. \gui{Query camera timing} reads
the \emph{actual} exposure and resulting frame rate from the IDS camera
(the sensor may clamp exposure below $1/\mathrm{fps}$ for readout).

\paragraph{Guided lab session (workflow buttons).}
\textbf{①~Background model (beam OFF)}: captures darks, builds the noise
model and keeps a dated ``before'' copy. \textbf{②~Start measuring (beam
ON)}: live neutron counting. \textbf{③~Camera health check (beam OFF)}:
fresh darks $\rightarrow$ ``after'' model $\rightarrow$ automatic radiation
damage report. \textbf{④~Background sanity}: classifies the beam-OFF
clusters $\rightarrow$ the ML false-positive floor per frame, to subtract
from any measured rate. The \emph{ambient} room rate (neighbouring
experiments) is measured by running ② with the beam OFF; ambient $=$ that
rate minus the ④ floor. \textbf{⑤~End-of-experiment report}: one click after
a run shows, in \gui{Results}, the four maps of the experiment that just
finished --- \emph{noise mean}, \emph{noise $\sigma$} (the sensor's
per-pixel background), \emph{arrival of all detected particles}, and
\emph{arrival of neutrons only} (positions of clusters classified
lithium/alpha) --- plus an automatic separability verdict (a red dialog
appears only if signal and noise have MERGED). Each ② run records into its
own timestamped folder (\texttt{data\_analysed/live\_<date>\_<time>}), so ⑤
always refers to exactly the experiment just finished, even with several
runs per day. ⑤ deliberately does \emph{not} produce the energy spectrum: a
normal 100--500-frame experiment has far too few heavy tracks for meaningful
peaks --- see the next paragraph.

\paragraph{Energy spectrum (office, large datasets).} The spectrum button
lives in \gui{Parameters} $\rightarrow$ \gui{Energy analysis}, directly above
a path field pointing at a LARGE classified CSV (e.g.\ the rare
15\,000-frame runs, or several sessions concatenated). Run it back at the
office, not in the live lab loop; the result appears in \gui{Results}.

\paragraph{Live telemetry.} The statistics line shows cumulative neutrons
($= \alpha + \mathrm{Li}$), per-channel counts and rates, neutrons/s,
neutrons/frame and dead-pixel statistics. The \textbf{throughput line} shows
\emph{Camera fps} (frames actually delivered), \emph{Analysed fps} (frames
actually processed), and \emph{Dropped}: \textcolor{green!50!black}{green} =
pipeline keeps up; \textcolor{orange}{orange} = queue filling, drops
imminent; \textcolor{red}{red} = frames are being lost and \emph{all
displayed rates under-count reality by that fraction} --- lower the camera
fps or relax the analysis. At the end of a live run the detector writes
\texttt{live\_signal\_arrival.tiff} (per-pixel capture counts; float32,
$\sim$80\,MB) and a PNG preview shown in \gui{Results} --- it reveals
\emph{where} captures land on the sensor (e.g.\ a partially covered region).

\subsection{Energy spectrum (Results tab)}
\gui{Energy Spectrum} runs \texttt{src/study\_energy.py} on
\texttt{detections\_classified.csv}: a histogram of \param{total\_charge}
for heavy tracks. Expected physics: two strong lines
($^{7}\mathrm{Li}$ 0.84\,MeV, $\alpha$ 1.47\,MeV, ratio 1.75) plus two weak
satellites (4\% branch: 1.01 and 1.78\,MeV). \textbf{Saturated pixels
destroy energy information at capture time} --- nothing can recover it in
post-processing; the tool prints the saturated fraction, and only spectra
with $\approx 0\%$ saturation (energy acquisition mode) are quantitative.
With \texttt{--calibrate} the two fitted peaks are anchored to the known
MeV lines, giving an absolute energy axis and the predicted satellite
positions.

\subsection{Camera damage tester (radiation health)}
Give two background models --- one built \emph{before} and one \emph{after}
irradiation (browse buttons accept the binary \texttt{.yml}/\texttt{.bgm}
files; the workflow's ①/③ buttons fill these in for you) --- and \gui{Run
Degradation Analysis} reports the dark-current shift, hot-pixel counts on
two independent signals, and a lifetime extrapolation.

Two hot-pixel signals are reported, deliberately: a raw \textbf{brightness
threshold} on the mean dark current (fast, but sensitive to how many dark
frames each model was built from --- the ``after'' model is always built
fresh, with no memory of the ``before'' one, so a small ``after'' sample
gives a noisier mean and can flip a borderline pixel across the threshold
either way by chance), and the \textbf{dead-pixel mask} each
\texttt{background\_model} run already builds internally from $\geq$95\%
saturated/zero behaviour over its \emph{entire} warmup, which cannot be
tripped by a single noisy frame. Radiation damage is one-directional: a
truly dead pixel cannot become alive again. If the dead-pixel count
\emph{drops} between ``before'' and ``after'', that is not the sensor
healing --- it is the measurement noise floor of that particular
before/after pair, and the tool says so explicitly. It also warns whenever
the two models were built from very different frame counts (ratio
$<0.5$ or $>2$), since that is exactly the condition that makes the
brightness-threshold signal unreliable. Build ``after'' from as many fresh
dark frames as time allows, ideally comparable to ``before''.

\section{Data products}
\begin{longtable}{p{5.2cm}p{9.5cm}}
\toprule
\textbf{File} & \textbf{Content} \\
\midrule
\texttt{background\_model.yml} & Binary noise model (mean, variance, counts, dead mask). \\
\texttt{detections.csv} & One row per detected cluster, 20 features. \\
\texttt{detections\_classified.csv} & Same + predicted label and confidence. \\
\texttt{data\_analysed/\ldots} & Compact archive: 5 raw sample frames,
sparse masked PNGs ($\sim$1000$\times$ smaller than raw --- only pixels of
valid clusters are kept), lossless float TIFFs of mean/$\sigma$/arrival.
Written by the offline steps and, when enabled, during live
(\texttt{data\_analysed/live\_<date>}) on a dedicated writer thread. \\
\texttt{live\_recording/\ldots} & Raw \texttt{.bmp} archive of every live
frame, when raw recording is enabled. \\
\texttt{live\_signal\_arrival.tiff/.png} & Where captures landed during the last live run. \\
\texttt{energy\_histogram.png} & Deposited-charge spectrum of heavy tracks. \\
\bottomrule
\end{longtable}

\section{Maintenance --- modifying and redeploying the code}
\begin{itemize}
  \item \textbf{Python change} (\texttt{PipelineController.py} or
        \texttt{Code/src/*.py}): nothing to build --- restart the GUI.
  \item \textbf{C++ change} (\texttt{*.cpp}/\texttt{*.h}): recompile before
        the change takes effect:
        \texttt{cmake --build build --config Release} inside \texttt{Code}
        (or simply re-run \texttt{WindowsLauncher.py}, which rebuilds then
        launches). The GUI runs whatever \texttt{build/detector(.exe)}
        currently is --- an old binary silently ignores new features.
  \item \textbf{New Python dependency}: re-run \texttt{PrepareSetupLab.py}
        on the online machine (it re-scans imports with pipreqs, sanitises
        the list, re-downloads wheels) and copy the new
        \texttt{Code/wheels} over.
  \item \textbf{Redeploying to the lab}: re-run \texttt{PrepareSetupLab.py}
        $\rightarrow$ fresh \texttt{ExportLab/} $\rightarrow$ USB. For a
        small patch, copying just the changed \texttt{.py} files is enough
        (plus a rebuild if C++ changed). Never copy \texttt{venv/} or
        \texttt{build/} between machines.
  \item \textbf{Changed gain / camera settings}: rebuild the background
        model (①) at the new operating point, and if the saturation regime
        changed, re-label and retrain the classifier (see \emph{Matched
        operating point}).
\end{itemize}

\section{Troubleshooting}
\begin{itemize}
  \item \emph{Neutron counts are zero}: retrain/re-export the classifier
        (\gui{Classify}) so \texttt{model\_bundle.onnx} exists; check the
        Console tab for engine errors.
  \item \emph{Camera fps lower than requested}: the number shown is what the
        camera \emph{actually delivers} (exposure time caps it at
        $1/\mathrm{exposure}$; disk decoding caps the simulated camera).
  \item \emph{Dropped frames (red)}: the analysis cannot keep up --- reduce
        camera fps, or raise \param{Min cluster size/charge}.
  \item \emph{``Corrupted model file''}: delete
        \texttt{background\_model.yml} and rerun \gui{Background Model}.
  \item \emph{Arrival heatmap looks binary (0/255)}: normal at low
        statistics --- each pixel has seen at most one capture; structure
        appears over long runs.
\end{itemize}

\end{document}
